\documentclass[12pt,a4paper]{article}

% ============================================================
% PACKAGES
% ============================================================
\usepackage[utf8]{inputenc}
\usepackage[T1]{fontenc}
\usepackage[french]{babel}
\usepackage{geometry}
\usepackage{graphicx}
\usepackage{amsmath}
\usepackage{amssymb}
\usepackage{array}
\usepackage{booktabs}
\usepackage{threeparttable}
\usepackage{longtable}
\usepackage{listings}
\usepackage{xcolor}
\usepackage{hyperref}
\usepackage{fancyhdr}
\usepackage{titlesec}
\usepackage{caption}
\usepackage{subcaption}
\usepackage{float}
\usepackage{enumitem}
\usepackage{tcolorbox}


% ============================================================
% MISE EN PAGE
% ============================================================
\geometry{
    left=2.5cm, right=2.5cm,
    top=2.5cm,  bottom=2.5cm
}

% En-tete et pied de page
\pagestyle{fancy}
\fancyhf{}
\rhead{\small Projet LMI ASD 2026}
\lhead{\small AIDA SENE --- COURA NDIONE}
\cfoot{\thepage}
\renewcommand{\headrulewidth}{0.4pt}

% Couleurs
\definecolor{bleu}{RGB}{26,58,92}
\definecolor{vert}{RGB}{29,158,117}
\definecolor{gris}{RGB}{245,245,245}
\definecolor{marron}{RGB}{168,45,45}
\definecolor{codecouleur}{RGB}{40,40,40}

% Style des sections
\titleformat{\section}
  {\Large\bfseries\color{bleu}}
  {\thesection}{1em}{}[\titlerule]

\titleformat{\subsection}
  {\large\bfseries\color{vert}}
  {\thesubsection}{1em}{}

\titleformat{\subsubsection}
  {\normalsize\bfseries\color{bleu}}
  {\thesubsubsection}{1em}{}

% Style du code C
\lstdefinestyle{styleC}{
    language=C,
    basicstyle=\ttfamily\small,
    keywordstyle=\bfseries\color{bleu},
    commentstyle=\itshape\color{gray},
    stringstyle=\color{vert},
    numberstyle=\tiny\color{gray},
    numbers=left,
    stepnumber=1,
    numbersep=8pt,
    backgroundcolor=\color{gris},
    frame=single,
    framerule=0.4pt,
    rulecolor=\color{bleu!30},
    breaklines=true,
    captionpos=b,
    tabsize=4,
    showstringspaces=false,
}
\lstset{style=styleC}

% Boite de remarque
\tcbuselibrary{skins}
\newtcolorbox{remarque}{
    colback=bleu!5,
    colframe=bleu,
    fonttitle=\bfseries,
    title=Remarque,
    boxrule=0.4pt
}

% ============================================================
% DOCUMENT
% ============================================================
\begin{document}

% ============================================================
% PAGE DE TITRE
% ============================================================
\begin{titlepage}
    \centering
    \vspace*{1cm}

    {\large\bfseries\color{bleu}
    Republique du Senegal \\
    Universite Iba Der Thiam de Thies \\
    UFR Sciences et Technologies \\
    Departement de Mathematiques et Informatique}

    \vspace{1.5cm}
    \hrule height 1pt
    \vspace{0.5cm}

    {\Huge\bfseries\color{bleu}
    Conception, Implementation et \\[0.3cm]
    Evaluation Comparative de \\[0.3cm]
    Structures de Donnees}

    \vspace{0.4cm}
    {\large\color{vert}
    Systeme de Gestion de Scolarite}

    \vspace{0.5cm}
    \hrule height 1pt
    \vspace{1.5cm}

    {\large
    \begin{tabular}{ll}
        \textbf{Filiere}     & Licence 2 Mathematiques--Informatique \\[6pt]
        \textbf{Module}      & Algorithmique et Structures de Donnees \\[6pt]
        \textbf{Enseignant}  & Dr Abdoulaye DIALLO \\[6pt]
        \textbf{Binome}      & AIDA SENE --- COURA DIONE \\[6pt]
        \textbf{Annee}       & 2025--2026 \\
    \end{tabular}}

    \vfill
    {\small\color{gray} Rapport redige en \LaTeX{} --- Annee universitaire 2025--2026}
\end{titlepage}

% ============================================================
% TABLE DES MATIERES
% ============================================================
\tableofcontents
\newpage

% ============================================================
% RESUME
% ============================================================
\begin{abstract}
Ce rapport presente la conception, l'implementation et l'evaluation
comparative de trois structures de donnees hybrides appliquees a un
systeme de gestion de scolarite : le tableau statique, le tableau
dynamique de pointeurs avec redimensionnement geometrique, et la liste
chainee avec sentinelle. Chaque structure supporte un ensemble
d'operations : insertion, recherche, suppression, tri (insertion et
rapide), agregations statistiques et persistance binaire. Une analyse
theorique rigoureuse de la complexite asymptotique est confrontee a
une etude experimentale reproductible portant sur des jeux de donnees
allant de $n = 100$ a $n = 100\,000$ enregistrements. L'etude
experimentale se concentre sur trois operations : la recherche, la
suppression et le tri (par insertion et rapide). Les resultats
confirment les predictions theoriques et mettent en evidence les
compromis temps/espace propres a chaque structure.
\end{abstract}

\newpage

% ============================================================
% 1. INTRODUCTION
% ============================================================
\section{Introduction}

\subsection{Contexte et motivation}

Dans tout systeme informatique realiste, le choix d'une representation
interne des donnees conditionne directement la performance des operations
courantes. Une decision de modelisation prise en debut de projet peut
multiplier par un facteur $100$ ou plus le temps d'execution lorsque
la volumetrie augmente.

Ce projet s'inscrit dans ce contexte : il place le developpeur dans
la posture d'un concepteur logiciel devant manipuler une base
d'enregistrements a plusieurs champs, organisee a l'aide de plusieurs
structures de donnees concurrentes. L'objectif est double : coder
correctement, mais aussi \textbf{mesurer, comparer et expliquer}
le comportement de chacune des structures retenues.

\subsection{Domaine applicatif}

Le domaine choisi est la \textbf{gestion de scolarite} : gestion des
etudiants, de leurs notes, inscriptions et informations personnelles.
Ce domaine est representatif des systemes d'information reels et
justifie l'utilisation de structures de donnees complexes reliees
par des pointeurs.

\subsection{Objectifs du projet}

\begin{enumerate}
    \item Concevoir un modele de donnees coherent avec des types
          d'enregistrements lies par pointeurs
    \item Implementer trois structures de donnees distinctes en C
    \item Effectuer une analyse de complexite asymptotique rigoureuse
    \item Mettre en place un protocole experimental reproductible
    \item Produire des courbes de complexite scientifiquement exploitables
\end{enumerate}

% ============================================================
% 2. MODELISATION DES DONNEES
% ============================================================
\section{Modelisation des Donnees}

\subsection{Choix de conception}

Le type principal manipule est \texttt{Etudiant}, accompagne de quatre
types secondaires : \texttt{Note}, \texttt{Inscription}, \texttt{Cours}
et \texttt{Date}. Cette organisation repond aux exigences du cahier
des charges : au moins deux types d'enregistrements lies par pointeurs.

\subsection{Structures de donnees definies}

\subsubsection{Type Date}

\begin{lstlisting}[caption={Structure Date}]
typedef struct {
    int jour;
    int mois;
    int annee;
} Date;
\end{lstlisting}

\subsubsection{Type Note}

\begin{lstlisting}[caption={Structure Note}]
typedef struct {
    int   matricule;
    float valeur;
    char  matiere[50];
    char *appreciation;    /* longueur variable -> malloc */
    Date  dateDeliberation;
    Cours *cours;          /* pointeur vers un Cours */
    int   tab[10];
} Note;
\end{lstlisting}

\subsubsection{Type Inscription}

\begin{lstlisting}[caption={Structure Inscription}]
typedef struct {
    int    matricule;
    float  montant;
    char   etat[30];
    char  *filliere;       /* longueur variable -> malloc */
    Date   dateInscription;
    Niveau *niveau;
    int    tab[20];
} Inscription;
\end{lstlisting}

\subsubsection{Type Etudiant (type principal)}

\begin{lstlisting}[caption={Structure Etudiant}]
typedef struct {
    int          matricule;
    float        moyenne;
    char         nom[40];
    char        *prenom;        /* longueur variable -> malloc */
    Date         dateNaissance;
    Note        *tabnote[10];   /* tableau de pointeurs vers Note */
    Inscription *inscription;   /* pointeur vers Inscription */
} Etudiant;
\end{lstlisting}

\subsection{Justification des choix}

Le champ \texttt{prenom} est declare comme \texttt{char*} (pointeur)
plutot que comme tableau fixe \texttt{char[n]}, afin de satisfaire
l'exigence d'un champ de longueur variable alloue dynamiquement avec
\texttt{malloc}. La memoire est allouee exactement a la taille necessaire :

\[
\text{taille} = \text{strlen(prenom)} + 1 \quad \text{(+1 pour le caractere nul)}
\]

Le tableau \texttt{tabnote[10]} constitue le champ tableau interne
non trivial exige, et le pointeur \texttt{inscription} assure la
relation $1$-$N$ entre les types d'enregistrements.

% ============================================================
% 3. STRUCTURES DE DONNEES IMPLEMENTEES
% ============================================================
\section{Structures de Donnees Implementees}

\subsection{Tableau statique}

\subsubsection{Principe}

Le tableau statique alloue un bloc continu de memoire a la compilation,
de capacite fixe \texttt{TAILLEMAX = 100\,000}. Les elements sont des
\texttt{Etudiant} (pas des pointeurs) stockes directement dans le tableau.

\subsubsection{Structure}

\begin{lstlisting}[caption={Tableau statique}]
Etudiant tab[TAILLEMAX];
int n = 0;  /* nombre d'elements actuellement stockes */
\end{lstlisting}

\subsubsection{Operations principales}

L'insertion maintient le tableau trie par matricule. La suppression
decale les elements suivants d'une position vers la gauche.

\subsection{Tableau dynamique de pointeurs}

\subsubsection{Principe}

Le tableau dynamique stocke des \textbf{pointeurs} vers des
\texttt{Etudiant}, et se redimensionne automatiquement par un
facteur de croissance de $\times 2$ via \texttt{realloc}.

\subsubsection{Structure}

\begin{lstlisting}[caption={Structure TableauDyn}]
typedef struct {
    Etudiant **tab;   /* tableau de pointeurs vers Etudiant */
    int taille;       /* nombre d'etudiants stockes         */
    int capacite;     /* capacite actuelle (avant realloc)  */
} TableauDyn;
\end{lstlisting}

\subsubsection{Redimensionnement geometrique}

Lorsque \texttt{taille == capacite}, la capacite est doublee :

\begin{lstlisting}[caption={Redimensionnement par realloc}]
int nouvelle_capacite = t->capacite * 2;
Etudiant **nouveau = realloc(t->tab,
                     nouvelle_capacite * sizeof(Etudiant *));
t->tab      = nouveau;
t->capacite = nouvelle_capacite;
\end{lstlisting}

La sequence des capacites est : $4 \to 8 \to 16 \to 32 \to \ldots$

\subsection{Liste chainee}

\subsubsection{Principe}

La liste chainee utilise des \textbf{noeuds} relies entre eux par un
pointeur \texttt{suiv}. Une \textbf{sentinelle} de debut (\texttt{tete})
elimine les cas particuliers d'insertion et de suppression en tete.

\subsubsection{Structures}

\begin{lstlisting}[caption={Structures Noeud et ListeC}]
typedef struct Noeud {
    Etudiant     *etudiant;
    struct Noeud *suiv;
} Noeud;

typedef struct {
    Noeud *tete;    /* sentinelle de debut */
    Noeud *queue;   /* dernier element     */
    int    taille;
} ListeC;
\end{lstlisting}

\subsubsection{Insertion en tete en O(1)}

Grace a la sentinelle, l'insertion en tete ne necessite aucun parcours :

\begin{lstlisting}[caption={Insertion en tete O(1)}]
Noeud *ancien_premier = l->tete->suiv;
nouveau->suiv  = ancien_premier;
l->tete->suiv  = nouveau;
l->taille++;
\end{lstlisting}

% ============================================================
% 4. OPERATIONS COMMUNES
% ============================================================
\section{Operations Communes}

\subsection{Recherche}

La recherche par cle primaire (matricule) est implementee pour chaque
structure par parcours sequentiel, et retourne un pointeur vers
l'etudiant ou \texttt{NULL} s'il n'est pas trouve.

\subsection{Tri}

Deux algorithmes de tri sont implementes :

\begin{itemize}
    \item \textbf{Tri par insertion} : trie par moyenne croissante.
          Performant sur les donnees presque triees.
    \item \textbf{Tri rapide (QuickSort)} : trie par moyenne croissante.
          Performant en moyenne sur les donnees aleatoires.
\end{itemize}

\subsection{Agregations statistiques}

Les agregations suivantes sont calculees sur le champ \texttt{moyenne} :

\[
\bar{x} = \frac{1}{n}\sum_{i=1}^{n} x_i \qquad
\sigma = \sqrt{\frac{1}{n}\sum_{i=1}^{n}(x_i - \bar{x})^2}
\]

\subsection{Persistance binaire}

La serialisation ecrit les champs fixes avec \texttt{fwrite},
et precede les champs de longueur variable de leur taille :

\begin{lstlisting}[caption={Serialisation du prenom}]
int len = strlen(e->prenom) + 1;
fwrite(&len,      sizeof(int),  1,   f);
fwrite(e->prenom, sizeof(char), len, f);
\end{lstlisting}

% ============================================================
% 5. ETUDE DE COMPLEXITE
% ============================================================
\section{Etude de Complexite}

\subsection{Analyse theorique}

\subsubsection{Tableau recapitulatif}

\begin{table}[H]
\centering
\caption{Complexite temporelle des operations par structure}
\begin{threeparttable}
\begin{tabular}{lccc}
\toprule
\textbf{Operation} & \textbf{Tab. statique} & \textbf{Tab. dynamique} & \textbf{Liste chainee} \\
\midrule
Recherche cle      & $O(n)/\Omega(1)$ & $O(n)/\Omega(1)$ & $O(n)/\Omega(1)$ \\
Suppression        & $\Theta(n)$   & $\Theta(n)$         & $O(n)/\Omega(1)$ \\
Tri insertion      & $O(n^2)/\Omega(n)$ & $O(n^2)/\Omega(n)$ & $O(n^2)/\Omega(n)$ \\
Tri rapide         & $O(n^2)/\Omega(n\log n)$ & $O(n^2)/\Omega(n\log n)$ & $O(n^2)/\Omega(n\log n)$ \\
\bottomrule
\end{tabular}
\end{threeparttable}
\end{table}

\subsubsection{Complexite spatiale}

\begin{table}[H]
\centering
\caption{Complexite spatiale par structure}
\begin{tabular}{lll}
\toprule
\textbf{Structure} & \textbf{Espace} & \textbf{Detail} \\
\midrule
Tableau statique   & $\Theta(C)$       & $C$ fixe a la compilation \\
Tableau dynamique  & $\Theta(n)$ a $\Theta(2n)$ & Au pire moitie vide apres realloc \\
Liste chainee      & $\Theta(n) + n$ pointeurs & Surcout constant par noeud \\
\bottomrule
\end{tabular}
\end{table}

\subsection{Justifications mathematiques}

\subsubsection{Tri par insertion}

La boucle externe itere de $i = 1$ a $n-1$. Pour chaque $i$,
la boucle interne effectue au plus $i$ comparaisons.
Le nombre total de comparaisons est :

\[
\sum_{i=1}^{n-1} i = \frac{n(n-1)}{2} = O(n^2)
\]

En meilleur cas (donnees deja triees), aucun decalage n'est necessaire :
complexite $\Omega(n)$.

\subsubsection{Tri rapide}

L'equation de recurrence en cas moyen est :
\[
T(n) = 2T\!\left(\frac{n}{2}\right) + O(n)
\]

Par le \textbf{theoreme maitre} (cas 2) avec $a=2$, $b=2$, $f(n)=n$ :
\[
n^{\log_2 2} = n^1 = n = f(n) \implies T(n) = \Theta(n \log n)
\]

Le pire cas $T(n) = O(n^2)$ survient lorsque la partition est
systematiquement desequilibree, c'est-a-dire lorsque le pivot choisi
tombe a une extremite de l'intervalle a chaque appel recursif.
Le choix du pivot au milieu de l'intervalle, utilise dans notre
implementation, limite fortement ce risque et rapproche le comportement
observe du cas moyen $\Theta(n \log n)$.

\subsection{Etude experimentale}

\subsubsection{Protocole}

\begin{itemize}
    \item \textbf{Tailles} : $n \in \{100,\ 1\,000,\ 10\,000,\ 100\,000\}$
    \item \textbf{Operations mesurees} : recherche, suppression, tri
          (par insertion et tri rapide)
    \item \textbf{Repetitions} : 10 mesures par configuration
    \item \textbf{Horloge} : \texttt{clock()} (bibliotheque standard C,
          precision de l'ordre de la milliseconde selon le systeme)
    \item \textbf{Resultats} : moyenne et ecart-type calcules
\end{itemize}

\begin{remarque}
Pour les petites tailles ($n=100$), la resolution de \texttt{clock()}
peut etre insuffisante pour distinguer des temps tres courts, ce qui
explique des mesures proches de zero. Une horloge a plus haute
resolution (\texttt{clock\_gettime(CLOCK\_MONOTONIC, ...)} sous
Linux/POSIX, ou \texttt{QueryPerformanceCounter} sous Windows)
donnerait une precision superieure et pourra etre envisagee dans
une version future du benchmark.
\end{remarque}

\subsubsection{Conditions de mesure}

\begin{table}[H]
\centering
\caption{Environnement de mesure}
\begin{tabular}{ll}
\toprule
\textbf{Parametre} & \textbf{Valeur} \\
\midrule
Systeme d'exploitation & Windows 10 \\
Memoire RAM            & 8 Go \\
Processeur            & Intel Core i5-1035G1 @ 1.00GHz \\
Compilateur            & GCC 10.3.0 TDM64 (-Wall -Wextra) \\
Langage                & C (standard C99) \\
\bottomrule
\end{tabular}
\end{table}

\subsubsection{Courbes de complexite}

Conformement au protocole experimental (section 5.3 du cahier des
charges), les courbes suivantes sont presentees : pour chaque
operation (recherche, suppression, tri par insertion, tri rapide),
une figure regroupant l'echelle lineaire (a gauche) et l'echelle
logarithmique (a droite), toutes structures superposees, ainsi
qu'une courbe de comparaison au modele theorique pour le tri
rapide. Un diagramme supplementaire presente le facteur de
ralentissement entre structures.

\paragraph{Temps d'execution par operation, toutes structures superposees}

\begin{figure}[H]
    \centering
    \includegraphics[width=\textwidth]{courbe_recherche.png}
    \caption{Temps de recherche en fonction de $n$ --- lineaire (gauche) et logarithmique (droite)}
\end{figure}

\begin{figure}[H]
    \centering
    \includegraphics[width=\textwidth]{courbe_suppression.png}
    \caption{Temps de suppression en fonction de $n$ --- lineaire (gauche) et logarithmique (droite)}
\end{figure}

\begin{figure}[H]
    \centering
    \includegraphics[width=\textwidth]{courbe_tri_insertion.png}
    \caption{Temps de tri par insertion en fonction de $n$ --- lineaire (gauche) et logarithmique (droite)}
\end{figure}

\begin{figure}[H]
    \centering
    \includegraphics[width=\textwidth]{courbe_tri_rapide.png}
    \caption{Temps de tri rapide en fonction de $n$ --- lineaire (gauche) et logarithmique (droite)}
\end{figure}

\paragraph{Comparaison au modele theorique}

\begin{figure}[H]
    \centering
    \includegraphics[width=0.9\textwidth]{comparaison_tri_rapide.png}
    \caption{Temps de tri rapide mesure compare a la courbe theorique $n\log n$}
\end{figure}

\paragraph{Facteur de ralentissement entre structures}

\begin{figure}[H]
    \centering
    \includegraphics[width=0.9\textwidth]{histogramme_ralentissement.png}
    \caption{Facteur de ralentissement entre la meilleure et la pire structure, par operation}
\end{figure}


% ============================================================
% 6. PERSISTANCE
% ============================================================
\section{Persistance des Donnees}

La serialisation binaire permet de sauvegarder l'integralite
d'un jeu de donnees sur disque et de le recharger ulterieurement.

\subsection{Format binaire}

Le format du fichier \texttt{.bin} est le suivant :

\begin{enumerate}
    \item Un entier \texttt{int} indiquant le nombre $n$ d'etudiants
    \item Pour chaque etudiant :
    \begin{itemize}
        \item Champs fixes : \texttt{matricule}, \texttt{moyenne},
              \texttt{nom[40]}, \texttt{dateNaissance}
        \item Champ variable : longueur du prenom (\texttt{int}),
              puis le prenom (\texttt{char*})
    \end{itemize}
\end{enumerate}

\subsection{Complexite de la persistance}

La serialisation et la deserialisation ont toutes deux une
complexite $\Theta(n)$ car elles parcourent chaque enregistrement
exactement une fois.

% ============================================================
% 7. CONCLUSION
% ============================================================
\section{Conclusion}

\subsection{Bilan comparatif}

Ce projet a permis d'implementer et de comparer trois structures
de donnees sur un meme jeu d'operations. Le tableau recapitulatif
des complexites theoriques a ete confirme par les mesures experimentales.

Les principaux enseignements sont :

\begin{itemize}
    \item Le \textbf{tableau statique} offre l'acces direct en $O(1)$
          mais souffre d'une capacite fixe et d'insertions couteuses
          lorsqu'il faut maintenir le tri.

    \item Le \textbf{tableau dynamique} equilibre flexibilite et
          performance grace au redimensionnement en $O(1)$ amorti,
          au prix d'une consommation memoire pouvant atteindre $2n$.

    \item La \textbf{liste chainee} excelle pour les insertions et
          suppressions en tete ($O(1)$), mais perd l'acces direct
          par index ($\Theta(n)$).
\end{itemize}

\subsection{Perspectives}

Des structures plus avancees comme les arbres binaires de recherche
equilibres (AVL, rouge-noir) ou les tables de hachage permettraient
d'atteindre une complexite $O(\log n)$ pour la recherche, au prix
d'une implementation plus complexe.

% ============================================================
% REFERENCES BIBLIOGRAPHIQUES
% ============================================================
\begin{thebibliography}{9}

\bibitem{cormen}
Cormen T. H., Leiserson C. E., Rivest R. L., Stein C.,
\textit{Introduction to Algorithms}, 4\textsuperscript{e} edition,
MIT Press, 2022.

\bibitem{sedgewick}
Sedgewick R., Wayne K.,
\textit{Algorithms}, 4\textsuperscript{e} edition,
Addison-Wesley, 2011.

\bibitem{knuth}
Knuth D. E.,
\textit{The Art of Computer Programming, vol. 3 : Sorting and Searching},
2\textsuperscript{e} edition, Addison-Wesley.

\bibitem{aho}
Aho A. V., Hopcroft J. E., Ullman J. D.,
\textit{Data Structures and Algorithms},
Addison-Wesley.

\bibitem{kr}
Kernighan B. W., Ritchie D. M.,
\textit{The C Programming Language},
2\textsuperscript{e} edition, Prentice Hall.

\bibitem{beauquier}
Beauquier D., Berstel J., Chretienne P.,
\textit{Elements d'algorithmique},
Masson.

\end{thebibliography}

\section*{Depot Git}
Le code source complet du projet est disponible publiquement a
l'adresse suivante :
\begin{center}
\url{https://github.com/TON_USERNAME/SENE-DIONE_SD}
\end{center}

\end{document}